\documentclass[11pt]{article}

    \usepackage[breakable]{tcolorbox}
    \usepackage{parskip} % Stop auto-indenting (to mimic markdown behaviour)
    

    % Basic figure setup, for now with no caption control since it's done
    % automatically by Pandoc (which extracts ![](path) syntax from Markdown).
    \usepackage{graphicx}
    % Keep aspect ratio if custom image width or height is specified
    \setkeys{Gin}{keepaspectratio}
    % Maintain compatibility with old templates. Remove in nbconvert 6.0
    \let\Oldincludegraphics\includegraphics
    % Ensure that by default, figures have no caption (until we provide a
    % proper Figure object with a Caption API and a way to capture that
    % in the conversion process - todo).
    \usepackage{caption}
    \DeclareCaptionFormat{nocaption}{}
    \captionsetup{format=nocaption,aboveskip=0pt,belowskip=0pt}

    \usepackage{float}
    \floatplacement{figure}{H} % forces figures to be placed at the correct location
    \usepackage{xcolor} % Allow colors to be defined
    \usepackage{enumerate} % Needed for markdown enumerations to work
    \usepackage{geometry} % Used to adjust the document margins
    \usepackage{amsmath} % Equations
    \usepackage{amssymb} % Equations
    \usepackage{textcomp} % defines textquotesingle
    % Hack from http://tex.stackexchange.com/a/47451/13684:
    \AtBeginDocument{%
        \def\PYZsq{\textquotesingle}% Upright quotes in Pygmentized code
    }
    \usepackage{upquote} % Upright quotes for verbatim code
    \usepackage{eurosym} % defines \euro

    \usepackage{iftex}
    \ifPDFTeX
        \usepackage[T1]{fontenc}
        \IfFileExists{alphabeta.sty}{
              \usepackage{alphabeta}
          }{
              \usepackage[mathletters]{ucs}
              \usepackage[utf8x]{inputenc}
          }
    \else
        \usepackage{fontspec}
        \usepackage{unicode-math}
    \fi

    \usepackage{fancyvrb} % verbatim replacement that allows latex
    \usepackage{grffile} % extends the file name processing of package graphics
                         % to support a larger range
    \makeatletter % fix for old versions of grffile with XeLaTeX
    \@ifpackagelater{grffile}{2019/11/01}
    {
      % Do nothing on new versions
    }
    {
      \def\Gread@@xetex#1{%
        \IfFileExists{"\Gin@base".bb}%
        {\Gread@eps{\Gin@base.bb}}%
        {\Gread@@xetex@aux#1}%
      }
    }
    \makeatother
    \usepackage[Export]{adjustbox} % Used to constrain images to a maximum size
    \adjustboxset{max size={0.9\linewidth}{0.9\paperheight}}

    % The hyperref package gives us a pdf with properly built
    % internal navigation ('pdf bookmarks' for the table of contents,
    % internal cross-reference links, web links for URLs, etc.)
    \usepackage{hyperref}
    % The default LaTeX title has an obnoxious amount of whitespace. By default,
    % titling removes some of it. It also provides customization options.
    \usepackage{titling}
    \usepackage{longtable} % longtable support required by pandoc >1.10
    \usepackage{booktabs}  % table support for pandoc > 1.12.2
    \usepackage{array}     % table support for pandoc >= 2.11.3
    \usepackage{calc}      % table minipage width calculation for pandoc >= 2.11.1
    \usepackage[inline]{enumitem} % IRkernel/repr support (it uses the enumerate* environment)
    \usepackage[normalem]{ulem} % ulem is needed to support strikethroughs (\sout)
                                % normalem makes italics be italics, not underlines
    \usepackage{soul}      % strikethrough (\st) support for pandoc >= 3.0.0
    \usepackage{mathrsfs}
    

    
    % Colors for the hyperref package
    \definecolor{urlcolor}{rgb}{0,.145,.698}
    \definecolor{linkcolor}{rgb}{.71,0.21,0.01}
    \definecolor{citecolor}{rgb}{.12,.54,.11}

    % ANSI colors
    \definecolor{ansi-black}{HTML}{3E424D}
    \definecolor{ansi-black-intense}{HTML}{282C36}
    \definecolor{ansi-red}{HTML}{E75C58}
    \definecolor{ansi-red-intense}{HTML}{B22B31}
    \definecolor{ansi-green}{HTML}{00A250}
    \definecolor{ansi-green-intense}{HTML}{007427}
    \definecolor{ansi-yellow}{HTML}{DDB62B}
    \definecolor{ansi-yellow-intense}{HTML}{B27D12}
    \definecolor{ansi-blue}{HTML}{208FFB}
    \definecolor{ansi-blue-intense}{HTML}{0065CA}
    \definecolor{ansi-magenta}{HTML}{D160C4}
    \definecolor{ansi-magenta-intense}{HTML}{A03196}
    \definecolor{ansi-cyan}{HTML}{60C6C8}
    \definecolor{ansi-cyan-intense}{HTML}{258F8F}
    \definecolor{ansi-white}{HTML}{C5C1B4}
    \definecolor{ansi-white-intense}{HTML}{A1A6B2}
    \definecolor{ansi-default-inverse-fg}{HTML}{FFFFFF}
    \definecolor{ansi-default-inverse-bg}{HTML}{000000}

    % common color for the border for error outputs.
    \definecolor{outerrorbackground}{HTML}{FFDFDF}

    % commands and environments needed by pandoc snippets
    % extracted from the output of `pandoc -s`
    \providecommand{\tightlist}{%
      \setlength{\itemsep}{0pt}\setlength{\parskip}{0pt}}
    \DefineVerbatimEnvironment{Highlighting}{Verbatim}{commandchars=\\\{\}}
    % Add ',fontsize=\small' for more characters per line
    \newenvironment{Shaded}{}{}
    \newcommand{\KeywordTok}[1]{\textcolor[rgb]{0.00,0.44,0.13}{\textbf{{#1}}}}
    \newcommand{\DataTypeTok}[1]{\textcolor[rgb]{0.56,0.13,0.00}{{#1}}}
    \newcommand{\DecValTok}[1]{\textcolor[rgb]{0.25,0.63,0.44}{{#1}}}
    \newcommand{\BaseNTok}[1]{\textcolor[rgb]{0.25,0.63,0.44}{{#1}}}
    \newcommand{\FloatTok}[1]{\textcolor[rgb]{0.25,0.63,0.44}{{#1}}}
    \newcommand{\CharTok}[1]{\textcolor[rgb]{0.25,0.44,0.63}{{#1}}}
    \newcommand{\StringTok}[1]{\textcolor[rgb]{0.25,0.44,0.63}{{#1}}}
    \newcommand{\CommentTok}[1]{\textcolor[rgb]{0.38,0.63,0.69}{\textit{{#1}}}}
    \newcommand{\OtherTok}[1]{\textcolor[rgb]{0.00,0.44,0.13}{{#1}}}
    \newcommand{\AlertTok}[1]{\textcolor[rgb]{1.00,0.00,0.00}{\textbf{{#1}}}}
    \newcommand{\FunctionTok}[1]{\textcolor[rgb]{0.02,0.16,0.49}{{#1}}}
    \newcommand{\RegionMarkerTok}[1]{{#1}}
    \newcommand{\ErrorTok}[1]{\textcolor[rgb]{1.00,0.00,0.00}{\textbf{{#1}}}}
    \newcommand{\NormalTok}[1]{{#1}}

    % Additional commands for more recent versions of Pandoc
    \newcommand{\ConstantTok}[1]{\textcolor[rgb]{0.53,0.00,0.00}{{#1}}}
    \newcommand{\SpecialCharTok}[1]{\textcolor[rgb]{0.25,0.44,0.63}{{#1}}}
    \newcommand{\VerbatimStringTok}[1]{\textcolor[rgb]{0.25,0.44,0.63}{{#1}}}
    \newcommand{\SpecialStringTok}[1]{\textcolor[rgb]{0.73,0.40,0.53}{{#1}}}
    \newcommand{\ImportTok}[1]{{#1}}
    \newcommand{\DocumentationTok}[1]{\textcolor[rgb]{0.73,0.13,0.13}{\textit{{#1}}}}
    \newcommand{\AnnotationTok}[1]{\textcolor[rgb]{0.38,0.63,0.69}{\textbf{\textit{{#1}}}}}
    \newcommand{\CommentVarTok}[1]{\textcolor[rgb]{0.38,0.63,0.69}{\textbf{\textit{{#1}}}}}
    \newcommand{\VariableTok}[1]{\textcolor[rgb]{0.10,0.09,0.49}{{#1}}}
    \newcommand{\ControlFlowTok}[1]{\textcolor[rgb]{0.00,0.44,0.13}{\textbf{{#1}}}}
    \newcommand{\OperatorTok}[1]{\textcolor[rgb]{0.40,0.40,0.40}{{#1}}}
    \newcommand{\BuiltInTok}[1]{{#1}}
    \newcommand{\ExtensionTok}[1]{{#1}}
    \newcommand{\PreprocessorTok}[1]{\textcolor[rgb]{0.74,0.48,0.00}{{#1}}}
    \newcommand{\AttributeTok}[1]{\textcolor[rgb]{0.49,0.56,0.16}{{#1}}}
    \newcommand{\InformationTok}[1]{\textcolor[rgb]{0.38,0.63,0.69}{\textbf{\textit{{#1}}}}}
    \newcommand{\WarningTok}[1]{\textcolor[rgb]{0.38,0.63,0.69}{\textbf{\textit{{#1}}}}}
    \makeatletter
    \newsavebox\pandoc@box
    \newcommand*\pandocbounded[1]{%
      \sbox\pandoc@box{#1}%
      % scaling factors for width and height
      \Gscale@div\@tempa\textheight{\dimexpr\ht\pandoc@box+\dp\pandoc@box\relax}%
      \Gscale@div\@tempb\linewidth{\wd\pandoc@box}%
      % select the smaller of both
      \ifdim\@tempb\p@<\@tempa\p@
        \let\@tempa\@tempb
      \fi
      % scaling accordingly (\@tempa < 1)
      \ifdim\@tempa\p@<\p@
        \scalebox{\@tempa}{\usebox\pandoc@box}%
      % scaling not needed, use as it is
      \else
        \usebox{\pandoc@box}%
      \fi
    }
    \makeatother

    % Define a nice break command that doesn't care if a line doesn't already
    % exist.
    \def\br{\hspace*{\fill} \\* }
    % Math Jax compatibility definitions
    \def\gt{>}
    \def\lt{<}
    \let\Oldtex\TeX
    \let\Oldlatex\LaTeX
    \renewcommand{\TeX}{\textrm{\Oldtex}}
    \renewcommand{\LaTeX}{\textrm{\Oldlatex}}
    % Document parameters
    % Document title
    \title{587-N14b}
    
    
    
    
    
    
    
% Pygments definitions
\makeatletter
\def\PY@reset{\let\PY@it=\relax \let\PY@bf=\relax%
    \let\PY@ul=\relax \let\PY@tc=\relax%
    \let\PY@bc=\relax \let\PY@ff=\relax}
\def\PY@tok#1{\csname PY@tok@#1\endcsname}
\def\PY@toks#1+{\ifx\relax#1\empty\else%
    \PY@tok{#1}\expandafter\PY@toks\fi}
\def\PY@do#1{\PY@bc{\PY@tc{\PY@ul{%
    \PY@it{\PY@bf{\PY@ff{#1}}}}}}}
\def\PY#1#2{\PY@reset\PY@toks#1+\relax+\PY@do{#2}}

\@namedef{PY@tok@w}{\def\PY@tc##1{\textcolor[rgb]{0.73,0.73,0.73}{##1}}}
\@namedef{PY@tok@c}{\let\PY@it=\textit\def\PY@tc##1{\textcolor[rgb]{0.24,0.48,0.48}{##1}}}
\@namedef{PY@tok@cp}{\def\PY@tc##1{\textcolor[rgb]{0.61,0.40,0.00}{##1}}}
\@namedef{PY@tok@k}{\let\PY@bf=\textbf\def\PY@tc##1{\textcolor[rgb]{0.00,0.50,0.00}{##1}}}
\@namedef{PY@tok@kp}{\def\PY@tc##1{\textcolor[rgb]{0.00,0.50,0.00}{##1}}}
\@namedef{PY@tok@kt}{\def\PY@tc##1{\textcolor[rgb]{0.69,0.00,0.25}{##1}}}
\@namedef{PY@tok@o}{\def\PY@tc##1{\textcolor[rgb]{0.40,0.40,0.40}{##1}}}
\@namedef{PY@tok@ow}{\let\PY@bf=\textbf\def\PY@tc##1{\textcolor[rgb]{0.67,0.13,1.00}{##1}}}
\@namedef{PY@tok@nb}{\def\PY@tc##1{\textcolor[rgb]{0.00,0.50,0.00}{##1}}}
\@namedef{PY@tok@nf}{\def\PY@tc##1{\textcolor[rgb]{0.00,0.00,1.00}{##1}}}
\@namedef{PY@tok@nc}{\let\PY@bf=\textbf\def\PY@tc##1{\textcolor[rgb]{0.00,0.00,1.00}{##1}}}
\@namedef{PY@tok@nn}{\let\PY@bf=\textbf\def\PY@tc##1{\textcolor[rgb]{0.00,0.00,1.00}{##1}}}
\@namedef{PY@tok@ne}{\let\PY@bf=\textbf\def\PY@tc##1{\textcolor[rgb]{0.80,0.25,0.22}{##1}}}
\@namedef{PY@tok@nv}{\def\PY@tc##1{\textcolor[rgb]{0.10,0.09,0.49}{##1}}}
\@namedef{PY@tok@no}{\def\PY@tc##1{\textcolor[rgb]{0.53,0.00,0.00}{##1}}}
\@namedef{PY@tok@nl}{\def\PY@tc##1{\textcolor[rgb]{0.46,0.46,0.00}{##1}}}
\@namedef{PY@tok@ni}{\let\PY@bf=\textbf\def\PY@tc##1{\textcolor[rgb]{0.44,0.44,0.44}{##1}}}
\@namedef{PY@tok@na}{\def\PY@tc##1{\textcolor[rgb]{0.41,0.47,0.13}{##1}}}
\@namedef{PY@tok@nt}{\let\PY@bf=\textbf\def\PY@tc##1{\textcolor[rgb]{0.00,0.50,0.00}{##1}}}
\@namedef{PY@tok@nd}{\def\PY@tc##1{\textcolor[rgb]{0.67,0.13,1.00}{##1}}}
\@namedef{PY@tok@s}{\def\PY@tc##1{\textcolor[rgb]{0.73,0.13,0.13}{##1}}}
\@namedef{PY@tok@sd}{\let\PY@it=\textit\def\PY@tc##1{\textcolor[rgb]{0.73,0.13,0.13}{##1}}}
\@namedef{PY@tok@si}{\let\PY@bf=\textbf\def\PY@tc##1{\textcolor[rgb]{0.64,0.35,0.47}{##1}}}
\@namedef{PY@tok@se}{\let\PY@bf=\textbf\def\PY@tc##1{\textcolor[rgb]{0.67,0.36,0.12}{##1}}}
\@namedef{PY@tok@sr}{\def\PY@tc##1{\textcolor[rgb]{0.64,0.35,0.47}{##1}}}
\@namedef{PY@tok@ss}{\def\PY@tc##1{\textcolor[rgb]{0.10,0.09,0.49}{##1}}}
\@namedef{PY@tok@sx}{\def\PY@tc##1{\textcolor[rgb]{0.00,0.50,0.00}{##1}}}
\@namedef{PY@tok@m}{\def\PY@tc##1{\textcolor[rgb]{0.40,0.40,0.40}{##1}}}
\@namedef{PY@tok@gh}{\let\PY@bf=\textbf\def\PY@tc##1{\textcolor[rgb]{0.00,0.00,0.50}{##1}}}
\@namedef{PY@tok@gu}{\let\PY@bf=\textbf\def\PY@tc##1{\textcolor[rgb]{0.50,0.00,0.50}{##1}}}
\@namedef{PY@tok@gd}{\def\PY@tc##1{\textcolor[rgb]{0.63,0.00,0.00}{##1}}}
\@namedef{PY@tok@gi}{\def\PY@tc##1{\textcolor[rgb]{0.00,0.52,0.00}{##1}}}
\@namedef{PY@tok@gr}{\def\PY@tc##1{\textcolor[rgb]{0.89,0.00,0.00}{##1}}}
\@namedef{PY@tok@ge}{\let\PY@it=\textit}
\@namedef{PY@tok@gs}{\let\PY@bf=\textbf}
\@namedef{PY@tok@ges}{\let\PY@bf=\textbf\let\PY@it=\textit}
\@namedef{PY@tok@gp}{\let\PY@bf=\textbf\def\PY@tc##1{\textcolor[rgb]{0.00,0.00,0.50}{##1}}}
\@namedef{PY@tok@go}{\def\PY@tc##1{\textcolor[rgb]{0.44,0.44,0.44}{##1}}}
\@namedef{PY@tok@gt}{\def\PY@tc##1{\textcolor[rgb]{0.00,0.27,0.87}{##1}}}
\@namedef{PY@tok@err}{\def\PY@bc##1{{\setlength{\fboxsep}{\string -\fboxrule}\fcolorbox[rgb]{1.00,0.00,0.00}{1,1,1}{\strut ##1}}}}
\@namedef{PY@tok@kc}{\let\PY@bf=\textbf\def\PY@tc##1{\textcolor[rgb]{0.00,0.50,0.00}{##1}}}
\@namedef{PY@tok@kd}{\let\PY@bf=\textbf\def\PY@tc##1{\textcolor[rgb]{0.00,0.50,0.00}{##1}}}
\@namedef{PY@tok@kn}{\let\PY@bf=\textbf\def\PY@tc##1{\textcolor[rgb]{0.00,0.50,0.00}{##1}}}
\@namedef{PY@tok@kr}{\let\PY@bf=\textbf\def\PY@tc##1{\textcolor[rgb]{0.00,0.50,0.00}{##1}}}
\@namedef{PY@tok@bp}{\def\PY@tc##1{\textcolor[rgb]{0.00,0.50,0.00}{##1}}}
\@namedef{PY@tok@fm}{\def\PY@tc##1{\textcolor[rgb]{0.00,0.00,1.00}{##1}}}
\@namedef{PY@tok@vc}{\def\PY@tc##1{\textcolor[rgb]{0.10,0.09,0.49}{##1}}}
\@namedef{PY@tok@vg}{\def\PY@tc##1{\textcolor[rgb]{0.10,0.09,0.49}{##1}}}
\@namedef{PY@tok@vi}{\def\PY@tc##1{\textcolor[rgb]{0.10,0.09,0.49}{##1}}}
\@namedef{PY@tok@vm}{\def\PY@tc##1{\textcolor[rgb]{0.10,0.09,0.49}{##1}}}
\@namedef{PY@tok@sa}{\def\PY@tc##1{\textcolor[rgb]{0.73,0.13,0.13}{##1}}}
\@namedef{PY@tok@sb}{\def\PY@tc##1{\textcolor[rgb]{0.73,0.13,0.13}{##1}}}
\@namedef{PY@tok@sc}{\def\PY@tc##1{\textcolor[rgb]{0.73,0.13,0.13}{##1}}}
\@namedef{PY@tok@dl}{\def\PY@tc##1{\textcolor[rgb]{0.73,0.13,0.13}{##1}}}
\@namedef{PY@tok@s2}{\def\PY@tc##1{\textcolor[rgb]{0.73,0.13,0.13}{##1}}}
\@namedef{PY@tok@sh}{\def\PY@tc##1{\textcolor[rgb]{0.73,0.13,0.13}{##1}}}
\@namedef{PY@tok@s1}{\def\PY@tc##1{\textcolor[rgb]{0.73,0.13,0.13}{##1}}}
\@namedef{PY@tok@mb}{\def\PY@tc##1{\textcolor[rgb]{0.40,0.40,0.40}{##1}}}
\@namedef{PY@tok@mf}{\def\PY@tc##1{\textcolor[rgb]{0.40,0.40,0.40}{##1}}}
\@namedef{PY@tok@mh}{\def\PY@tc##1{\textcolor[rgb]{0.40,0.40,0.40}{##1}}}
\@namedef{PY@tok@mi}{\def\PY@tc##1{\textcolor[rgb]{0.40,0.40,0.40}{##1}}}
\@namedef{PY@tok@il}{\def\PY@tc##1{\textcolor[rgb]{0.40,0.40,0.40}{##1}}}
\@namedef{PY@tok@mo}{\def\PY@tc##1{\textcolor[rgb]{0.40,0.40,0.40}{##1}}}
\@namedef{PY@tok@ch}{\let\PY@it=\textit\def\PY@tc##1{\textcolor[rgb]{0.24,0.48,0.48}{##1}}}
\@namedef{PY@tok@cm}{\let\PY@it=\textit\def\PY@tc##1{\textcolor[rgb]{0.24,0.48,0.48}{##1}}}
\@namedef{PY@tok@cpf}{\let\PY@it=\textit\def\PY@tc##1{\textcolor[rgb]{0.24,0.48,0.48}{##1}}}
\@namedef{PY@tok@c1}{\let\PY@it=\textit\def\PY@tc##1{\textcolor[rgb]{0.24,0.48,0.48}{##1}}}
\@namedef{PY@tok@cs}{\let\PY@it=\textit\def\PY@tc##1{\textcolor[rgb]{0.24,0.48,0.48}{##1}}}

\def\PYZbs{\char`\\}
\def\PYZus{\char`\_}
\def\PYZob{\char`\{}
\def\PYZcb{\char`\}}
\def\PYZca{\char`\^}
\def\PYZam{\char`\&}
\def\PYZlt{\char`\<}
\def\PYZgt{\char`\>}
\def\PYZsh{\char`\#}
\def\PYZpc{\char`\%}
\def\PYZdl{\char`\$}
\def\PYZhy{\char`\-}
\def\PYZsq{\char`\'}
\def\PYZdq{\char`\"}
\def\PYZti{\char`\~}
% for compatibility with earlier versions
\def\PYZat{@}
\def\PYZlb{[}
\def\PYZrb{]}
\makeatother


    % For linebreaks inside Verbatim environment from package fancyvrb.
    \makeatletter
        \newbox\Wrappedcontinuationbox
        \newbox\Wrappedvisiblespacebox
        \newcommand*\Wrappedvisiblespace {\textcolor{red}{\textvisiblespace}}
        \newcommand*\Wrappedcontinuationsymbol {\textcolor{red}{\llap{\tiny$\m@th\hookrightarrow$}}}
        \newcommand*\Wrappedcontinuationindent {3ex }
        \newcommand*\Wrappedafterbreak {\kern\Wrappedcontinuationindent\copy\Wrappedcontinuationbox}
        % Take advantage of the already applied Pygments mark-up to insert
        % potential linebreaks for TeX processing.
        %        {, <, #, %, $, ' and ": go to next line.
        %        _, }, ^, &, >, - and ~: stay at end of broken line.
        % Use of \textquotesingle for straight quote.
        \newcommand*\Wrappedbreaksatspecials {%
            \def\PYGZus{\discretionary{\char`\_}{\Wrappedafterbreak}{\char`\_}}%
            \def\PYGZob{\discretionary{}{\Wrappedafterbreak\char`\{}{\char`\{}}%
            \def\PYGZcb{\discretionary{\char`\}}{\Wrappedafterbreak}{\char`\}}}%
            \def\PYGZca{\discretionary{\char`\^}{\Wrappedafterbreak}{\char`\^}}%
            \def\PYGZam{\discretionary{\char`\&}{\Wrappedafterbreak}{\char`\&}}%
            \def\PYGZlt{\discretionary{}{\Wrappedafterbreak\char`\<}{\char`\<}}%
            \def\PYGZgt{\discretionary{\char`\>}{\Wrappedafterbreak}{\char`\>}}%
            \def\PYGZsh{\discretionary{}{\Wrappedafterbreak\char`\#}{\char`\#}}%
            \def\PYGZpc{\discretionary{}{\Wrappedafterbreak\char`\%}{\char`\%}}%
            \def\PYGZdl{\discretionary{}{\Wrappedafterbreak\char`\$}{\char`\$}}%
            \def\PYGZhy{\discretionary{\char`\-}{\Wrappedafterbreak}{\char`\-}}%
            \def\PYGZsq{\discretionary{}{\Wrappedafterbreak\textquotesingle}{\textquotesingle}}%
            \def\PYGZdq{\discretionary{}{\Wrappedafterbreak\char`\"}{\char`\"}}%
            \def\PYGZti{\discretionary{\char`\~}{\Wrappedafterbreak}{\char`\~}}%
        }
        % Some characters . , ; ? ! / are not pygmentized.
        % This macro makes them "active" and they will insert potential linebreaks
        \newcommand*\Wrappedbreaksatpunct {%
            \lccode`\~`\.\lowercase{\def~}{\discretionary{\hbox{\char`\.}}{\Wrappedafterbreak}{\hbox{\char`\.}}}%
            \lccode`\~`\,\lowercase{\def~}{\discretionary{\hbox{\char`\,}}{\Wrappedafterbreak}{\hbox{\char`\,}}}%
            \lccode`\~`\;\lowercase{\def~}{\discretionary{\hbox{\char`\;}}{\Wrappedafterbreak}{\hbox{\char`\;}}}%
            \lccode`\~`\:\lowercase{\def~}{\discretionary{\hbox{\char`\:}}{\Wrappedafterbreak}{\hbox{\char`\:}}}%
            \lccode`\~`\?\lowercase{\def~}{\discretionary{\hbox{\char`\?}}{\Wrappedafterbreak}{\hbox{\char`\?}}}%
            \lccode`\~`\!\lowercase{\def~}{\discretionary{\hbox{\char`\!}}{\Wrappedafterbreak}{\hbox{\char`\!}}}%
            \lccode`\~`\/\lowercase{\def~}{\discretionary{\hbox{\char`\/}}{\Wrappedafterbreak}{\hbox{\char`\/}}}%
            \catcode`\.\active
            \catcode`\,\active
            \catcode`\;\active
            \catcode`\:\active
            \catcode`\?\active
            \catcode`\!\active
            \catcode`\/\active
            \lccode`\~`\~
        }
    \makeatother

    \let\OriginalVerbatim=\Verbatim
    \makeatletter
    \renewcommand{\Verbatim}[1][1]{%
        %\parskip\z@skip
        \sbox\Wrappedcontinuationbox {\Wrappedcontinuationsymbol}%
        \sbox\Wrappedvisiblespacebox {\FV@SetupFont\Wrappedvisiblespace}%
        \def\FancyVerbFormatLine ##1{\hsize\linewidth
            \vtop{\raggedright\hyphenpenalty\z@\exhyphenpenalty\z@
                \doublehyphendemerits\z@\finalhyphendemerits\z@
                \strut ##1\strut}%
        }%
        % If the linebreak is at a space, the latter will be displayed as visible
        % space at end of first line, and a continuation symbol starts next line.
        % Stretch/shrink are however usually zero for typewriter font.
        \def\FV@Space {%
            \nobreak\hskip\z@ plus\fontdimen3\font minus\fontdimen4\font
            \discretionary{\copy\Wrappedvisiblespacebox}{\Wrappedafterbreak}
            {\kern\fontdimen2\font}%
        }%

        % Allow breaks at special characters using \PYG... macros.
        \Wrappedbreaksatspecials
        % Breaks at punctuation characters . , ; ? ! and / need catcode=\active
        \OriginalVerbatim[#1,codes*=\Wrappedbreaksatpunct]%
    }
    \makeatother

    % Exact colors from NB
    \definecolor{incolor}{HTML}{303F9F}
    \definecolor{outcolor}{HTML}{D84315}
    \definecolor{cellborder}{HTML}{CFCFCF}
    \definecolor{cellbackground}{HTML}{F7F7F7}

    % prompt
    \makeatletter
    \newcommand{\boxspacing}{\kern\kvtcb@left@rule\kern\kvtcb@boxsep}
    \makeatother
    \newcommand{\prompt}[4]{
        {\ttfamily\llap{{\color{#2}[#3]:\hspace{3pt}#4}}\vspace{-\baselineskip}}
    }
    

    
    % Prevent overflowing lines due to hard-to-break entities
    \sloppy
    % Setup hyperref package
    \hypersetup{
      breaklinks=true,  % so long urls are correctly broken across lines
      colorlinks=true,
      urlcolor=urlcolor,
      linkcolor=linkcolor,
      citecolor=citecolor,
      }
    % Slightly bigger margins than the latex defaults
    
    \geometry{verbose,tmargin=1in,bmargin=1in,lmargin=1in,rmargin=1in}
    
    

\begin{document}
    
    \maketitle
    
    

    
    \hypertarget{material-balances-ii}{%
\section{Material Balances II}\label{material-balances-ii}}

\href{https://www.youtube.com/watch?v=7vP2hFFV57E\&list=RD7vP2hFFV57E\&start_radio=1}{What
is Bond listening to right now?}

This lecture continues discussion of material balances on reacting
systems.

In
\href{https://jqbond.github.io/CEN587/Notebooks/587-N13.html}{Notebook
13}, we looked at a species balance on a reactor of arbitrary ``size''
and shape, where we considered that the ``size'' of the system might be
described in terms of volume, mass of catalyst, interfacial area, or
even number of catalytic ``active sites'' in the reactor. Our species
balance on an arbitrary reactor:

\begin{figure}
\centering
\includegraphics{MB1.png}
\caption{Conceptual Illustration of an Arbitrary Chemical Reactor}
\end{figure}

Would give the following equation:

\[\frac{dN_j}{dt} = F_{j,f} - F_j + G_j\]

In other words, \{accumulation\} = \{in\} - \{out\} + \{generation by
reactions\}. We finished
\href{https://jqbond.github.io/CEN587/Notebooks/587-N13.html\#the-rate-of-generation-or-consumption-by-chemical-reaction}{Notebook
13} with a brief presentation of \(G_j\). We willl pick this up here.

\(G_j\) is the \textbf{\emph{extensive}} rate of generation (or
consumption) of species \(j\) by chemical reaction. It represents how
quickly species \(j\) is either consumed or produced by chemical
reaction(s) that are occuring inside of our reactor in
\textbf{\emph{extensive}} units of \textbf{\emph{moles-per-time}}. It is
related to the rate of reaction, but it is not the same thing as the
rate of reaction:

\[G_j \neq r_i\]

\(G_j\) is a \textbf{\emph{species specific term}}. It is an
\textbf{\emph{extensive}} quantity that represents the amount of species
\(j\) that is produced or consumed by chemical reaction in extensive
units of moles-per-time. It will change depending on the size of the
system and the particulars of the reactor you are using. In contrast,
the rate of reaction, \(r_i\), is a \textbf{\emph{reaction specific
term}}. It represents the rate of reaction \(i\) that is occuring at the
conditions inside of the reactor you are considering. Importantly, in
Chemical Reaction Engineering, the rate of a chemical reaction, \(r_i\),
is always an \textbf{\emph{intensive}} quantity that is scale and/or
system independent. In terms of what controls the intensive rate of
reaction, \(r_i\):

\[r_i = f(T, P, \chi_j)\]

To expand that concept: once we decide on what reaction we're
considering and what (if any) catalyst it uses, \textbf{\emph{the rate
of that reaction, \(r_i\), is entirely determined by the temperature,
pressure, and composition of the reacting mixture}}. The rate of
reaction, \(r_i\), does not depend on stoichiometry, nor does it depend
on the type, size, or shape of a reactor. It might surprise you, but the
rate of reaction also has no fundamental dependence on important aspects
of reactor design, like mixing. The rate of reaction is always entirely
determined by temperature, pressure, and composition. In contrast,
stoichiometry, shape, size, type of reactor, extent of mixing, etc. will
all influence the extensive rate of generation by chemical reaction,
\(G_j\):

\[G_j = f(T, P, \chi_j, \text{reactor type, reactor size, mixing, etc.})\]

These two quantities are related, but they are not the same thing. Next,
we'll develop the relationship between the two.

    \hypertarget{a-conceptual-definition-of-a-reaction-rate}{%
\subsection{A conceptual definition of a reaction
rate}\label{a-conceptual-definition-of-a-reaction-rate}}

I have always liked the way that reaction rates are conceptually defined
in Boudart's textbook \emph{Kinetics of Chemical Processes}. They are
not defined in terms of a species, rather, they are defined in terms of
the amount of ``reaction'' that has occured. Namely reaction rates are
defined in terms of an extent of reaction, \(\varepsilon\), which we
know is an extensive quantity that describes how many ``moles of
reaction'' have occured. We want to transform that idea into a rate,
specifically, into a rate of how many moles of reaction have occurred.
At least from a dimensional standpoint, one could define an extensive
rate of reaction \(i\), \(\bar{r_i}\) as:

\[\bar{r_i} = \frac{\Delta \varepsilon_i}{\Delta t} \ [=] \ \frac{\text{moles of extent of reaction} \ i}{\text{time}}\]

Since the extent of reaction \(i\), \(\varepsilon_i\), is an extensive
variable, the above definition of \(\bar{r_i}\) is an extensive rate of
reaction, i.e., it varies with the scale of the system.

\hypertarget{intensive-rates-of-reaction-are-more-useful}{%
\subsection{Intensive rates of reaction are more
useful}\label{intensive-rates-of-reaction-are-more-useful}}

Generally, we prefer to work with \emph{intensive} reaction rates. The
reason for this is that, fundamentally, a reaction rate should be a
function of only \((T, P, x_j)\). It should not depend on the type of
reactor we are using or the size of that system. This ensures that when
I study kinetics in a small scale batch reactor, those results will
apply to the same reaction happening in a large fluidized bed at
ExxonMobil. Intensive reaction rates ensure the data are universal and
that they will apply for a given \((T, P, x_j)\) regardless of the
reactor. This is a much more useful description of reaction rate than an
extensive reaction rate that changes every time we change the reactor.
Accordingly, in chemical reactor design, we will always work with
intensive reaction rates that are normalized to a unit dimension of
``size'' or ``scale'' that describes our system. As we discussed in
\href{https://jqbond.github.io/CEN587/Notebooks/587-N13.html}{Notebook
13}, what constitutes an appropriate dimensional scaling depends on the
type of reaction. For homogeneous reactions that occur throughout the
space occupied by the reaction mixture, volume is a good description of
system size. We reach this conclusion because it is clear that
increasing the volume of fluid within the reactor (for a fixed
concentration of reactants) will increase the extensive quantity of
reaction that is occuring. In other cases, catalyst mass, interfacial
surface area, or even number of surface active sites may be appropriate
bases for normalizing reaction rates. With this in mind, we present a
few general conceptual definitions of intensive reaction rates using
different types of normalizations that may be relevant in different
types of reacting systems.

\hypertarget{homogeneous-reactions}{%
\subsubsection{Homogeneous Reactions}\label{homogeneous-reactions}}

For homogeneous reactions that occur throughout the space occupied by a
fluid phase, volume is a good basis for normalizing reaction rates;
hence, our convention for homogeneous systems is to report intensive
reaction rates normalized to unit volume:

\[r_i = \frac{\Delta \varepsilon_i}{\Delta t \ \Delta V} \ [=] \ \frac{\text{moles of extent of reaction} \ i}{\text{volume} \cdot \text{time}}\]

We will use the notation convention presented in Scott Fogler's
\emph{Elements of Chemical Reaction Engineering} wherein the rate of
reaction \(i\) reported per unit volume will have the symbol \(r_i\).

\hypertarget{heterogeneous-reactions}{%
\subsubsection{Heterogeneous Reactions}\label{heterogeneous-reactions}}

\hypertarget{catalytic-reactions-e.g.-in-a-packed-bed}{%
\paragraph{Catalytic reactions (e.g., in a packed
bed)}\label{catalytic-reactions-e.g.-in-a-packed-bed}}

In cases where the reaction occurs on a catalyst, reactor volume is a
much less relevant description of the reactor scale than the amount of
catalyst present in the reactor. In practice, for catalytic reactions,
we frequently will see intensive reaction rates reported per unit mass
of catalyst. This makes sense because the \textbf{\emph{extensive}}
productivity of a catalytic reactor scales with the quantity (mass) of
catalyst inside of the reactor much moreso than it does the volume of
space inside of the reactor. In this case, we define:

\[r_i^\prime = \frac{\Delta \varepsilon_i}{\Delta t \ \Delta W} \ [=] \ \frac{\text{moles of extent of reaction} \ i}{\text{mass of catalyst} \cdot \text{time}}\]

Note that Scott Fogler's textbook (\emph{Elements of Chemical Reaction
Engineering}) uses the symbol \(W\) for catalyst mass (catalyst
``weight''). This has become standard notation for packed bed and
fluidized bed reactors, so we'll use it here. In our course, we'll also
use Fogler's convention of using the symbol \(r_i^\prime\) to indicate
the mass-normalized rate of reaction \(i\).

\hypertarget{reactions-at-interfaces-between-phases}{%
\paragraph{Reactions at interfaces between
phases}\label{reactions-at-interfaces-between-phases}}

Sometimes, a reaction occurs only at an interface between two phases. In
these cases, the extensive ``scale'' of the reaction is best described
in terms of that interfacial area. When we have a reaction that occurs
over a 2D surface or 2D interface between phases, it is often convenient
for us to report reaction rates per unit area. Again, this is because
the interfacial area is far more important than the total volume
occupied in determining the overall, extensive performance of the
reactor. In these cases, we might see a rate definition like this one:

\[r_i^{\prime\prime} = \frac{\Delta \varepsilon_i}{\Delta t \ \Delta A} \ [=] \ \frac{\text{moles of extent of reaction} \ i}{\text{interfacial area} \cdot \text{time}}\]

We will extend the prime notation used by Fogler as well as Dumesic (in
\emph{The Microkinetics of Heterogeneous Catalysis}) and use the symbol
\(r_i^{\prime\prime}\) to represent the rate of reaction \(i\) per unit
area.

\hypertarget{reactions-at-specific-active-sites-on-a-surface}{%
\paragraph{Reactions at specific active sites on a
surface}\label{reactions-at-specific-active-sites-on-a-surface}}

In
\href{https://jqbond.github.io/CEN587/Notebooks/587-N13.html\#heterogenous-reactions-at-exposed-surface-atoms-active-sites}{Notebook
13}, we introduced the idea that catalysts are frequently comprised of
expensive metals like Pt or Pd. In those cases, we almost never use a
bulk sample of pure Pt or Pd. Instead, we disperse nanometer sized
particles of Pt or Pd onto a high surface area carrier like silica or
alumina. Frequently, the silica and alumina do nothing in terms of
catalysis--they are there as a host for Pt or Pd nanoparticles. In these
situations, the vast majority of catalyst mass or even catalyst surface
area in the reactor is inert. Only a few exposed surface atoms of Pt or
Pd are doing the catalysis, so the mass of catalyst and the surface area
of the interface between the catalyst and the bulk fluid become poor
descriptions of the size of the system. Instead, it is better for us to
think of the reactor ``size'' in terms of the number of Pt or Pd atoms
that are accessible in that system. This is a common convention in
fundamental catalysis, where we usually report rates of reaction per
active site:

\[r_i^{\prime\prime\prime} = \frac{\Delta \varepsilon_i}{\Delta t \ \Delta S} \ [=] \ \frac{\text{moles of extent of reaction} \ i}{\text{active site} \cdot \text{time}}\]

Usually, we quantity our active site in number of moles of active sites:

\[r_i^{\prime\prime\prime} = \frac{\Delta \varepsilon_i}{\Delta t \ \Delta S} \ [=] \ \frac{\text{moles of extent of reaction} \ i}{\text{moles of active site} \cdot \text{time}}\]

Which leads to a common (though not strictly rigorous) ``cancellation''
of moles to give a reaction rate in units of inverse time. This is
called the turnover frequency of a catalytic reaction.

\[r_i^{\prime\prime\prime} = \frac{\Delta \varepsilon_i}{\Delta t \ \Delta S} \ [=] \ \frac{\text{1}}{\text{time}}\]

The notation here is consistent with that presented in Dumesic's The
Microkinetics of Heterogeneous Catalysis, where
\(r_i^{\prime\prime\prime}\) represents the rate of reaction \(i\) per
unit active site. It can be reduced to units of inverse time, and is is
also known as the \textbf{\emph{turnover frequency}} of a catalytic
reaction.

There are many types of catalysts, and they have many types of ``active
sites.'' Above, we present the idea that a catalyst might be comprised
of a metal dispersed on an inert carrier and that the active sites are
accessible, surface metal atoms (not metal atoms buried in the bulk of a
3D particle). Some common examples of catalytic metals include Pt, Pd,
Ir, Ru, Ni, Cu, and Os. Sometimes, we use metal oxides or mixed metal
oxides as catalysts. In these cases, the active sites are often
coordinatively unsaturated Lewis acidic metal cations (electron
acceptors) like \(Al^{3+}\); Brønsted acidic protons, \(H^{+}\); basic
oxygen anions like \(O^{2-}\); or reducible metal cations like
\(V^{5+}\) or \(Ti^{4+}\). In all cases, the active site is likely a
rare species on the surface of the catalyst.

    \hypertarget{relating-production-rates-to-reaction-rates}{%
\subsection{Relating Production Rates to Reaction
Rates}\label{relating-production-rates-to-reaction-rates}}

For the rest of this lecture, we'll work with a homogeneous reaction and
so consider the development of intensive reaction rates that are
normalized per unit volume, i.e.,

\[r_i = \frac{\Delta \varepsilon_i}{\Delta t \ \Delta V} \ [=] \ \frac{\text{moles of extent of reaction} \ i}{\text{volume} \cdot \text{time}}\]

Consider a homogeneous reaction--let's call it Reaction 1:

\[2A \ (g) \longrightarrow 3B \ (g)\]

Under the conditions of our system, Reaction 1 is occuring with an
\textbf{\emph{intensive}} reaction rate of:

\[r_1 = 1 \ \frac{\text{moles of extent of reaction} \ 1}{\text{L} \cdot \text{min}}\]

Remember, we have defined reaction rates as properties of the reaction
(not of a species), and they represent the rate of change in extent of
reaction per unit time and size. In material balances, though, we are
always interested in the rate that \textbf{\emph{species}} are being
generated by chemical reaction more than we are interested in the rate
of the reaction itself. We can relate the intensive rate of reaction
\(i\), \(r_i\), to the intensive production rate of species \(j\),
\(R_j\), using stoichiometric coefficients. This is similar to what we
do with extents of reaction in completing a mole table for equilibrium
problems. For example, if we wanted to define the number of moles of
\(A\) and \(B\) as functions of extent of reaction 1, \(\varepsilon_1\),
we would write:

\begin{align}
    N_A = N_{A0} + \nu_A \varepsilon = N_{A0} - 2 \varepsilon_1 \\
    N_B = N_{B0} + \nu_B \varepsilon = N_{B0} + 3 \varepsilon_1 \\
\end{align}

From this, you can see how stoichiometric coefficients allow us to
relate a reaction specific property (like extent) to a species specific
property (number of moles of \(A\) consumed). We can use an anlogous
approach to relate reaction-specific reaction rates, \(r_i\), to
species-specific production rates, \(R_j\). For example, if we wanted to
calculate the rate at which species A is produced, \(R_A\), by reaction
1, we would consider that each time reaction 1 occurs, it consumes 2
moles of species A. \(R_A\) would therefore be defined as:

\[R_A = \nu_A r_1 = -2r_1 = -2 \ \frac{\text{moles \ of \ A}}{\text{L} \cdot \text{min}}\]

Similarly, the production rate of species B, \(R_B\), is given by:

\[R_B = \nu_B r_1 = 3r_1 = 3 \ \frac{\text{moles \ of \ B}}{\text{L} \cdot \text{min}}\]

Notice that these quantities, \(R_j\), are defined per unit size (volume
in this case), so these capital \(R\) production rates, \(R_j\) are
intensive quantities, just like reaction rates, \(r_i\). Intensive
production rates for homogeneous reactions will usually have the
following units:

\[R_j \ [=] \ \frac{\text{moles of j}}{\text{volume} \cdot \text{time}}\]

\hypertarget{generalizing-the-concept}{%
\subsubsection{Generalizing the
Concept}\label{generalizing-the-concept}}

Based on that example above, we can conclude that for any, generic
single reaction:

\[\nu_A A + \nu_B B \leftrightharpoons \nu_C C + \nu_D D\]

That is occuring with an intensive reaction rate, \(r_i\). We would
define the intensive production rate of a species participating in that
reaction as:

\[R_j = \nu_j r_i\]

Conversely, in some situations, we are able to estimate the rate of
reaction (which we can't measure directly) from the production rate of a
species (which we can measure directly):

\[r_i = \frac{R_j}{\nu_j}\]

    \hypertarget{multiple-reactions}{%
\subsubsection{Multiple Reactions}\label{multiple-reactions}}

There is no restriction that a species is only participating in one
reaction in a given reactor. In fact, the more realistic scenario is
that each species is actually participating in multiple reactions. This
means that each species is being consumed or produced by multiple
chemical reactions, and so when we define the production rate, \(R_j\),
we actually need to account for the rate of consumption from many
chemical reactions. We'll do this with an intuitive example and then
present a general equation. Consider partial and complete combustion
pathways that occur during methane oxidation:

\begin{align*}
    (1) \qquad CH_4 \ (g) + \frac{3}{2} O_2 \ (g) &\longrightarrow CO \ (g) + 2H_2O \ (g) \\
    \\
    (2) \qquad CH_4 \ (g) + 2 O_2 \ (g) &\longrightarrow CO_2 \ (g) + 2H_2O \ (g) \\
\end{align*}

These are different reactions, and their rates will have different
dependencies on temperature, pressure, and composition. In other words:

\[r_1(T,P,x_j) \neq r_2(T,P,x_j)\]

Because of this, we have to account for the rate of each reaction
separately when we quantify the production rate for each species.
Conceptually, it's pretty straightforward. Analogous to our solution of
multiple reaction equilibria using mole tables, we simply add up the
consumption or production rate for each species in each reaction. So for
methane \((CH_4)\), we find:

\[R_{CH_4} = -1r_1 - 1r_2\]

And we can do the same thing for every species participating in these
two reactions.

\begin{align}
    R_{O_2} &= -\frac{3}{2}r_1 - 2r_2\\
    R_{CO} &= +1r_1 + 0r_2\\
    R_{CO_2} &= +0r_1 + 1r_2\\
    R_{H_2O} &= +2r_1 + 2r_2\\
\end{align}

From that exercise, it makes sense that we can formally calculate an
intensive production rate for any species, \(j\), by summing up the
intensive production rate for that species in each reaction that is
occuring in the system. That results in the following definition of
intensive production rates:

\[R_j = \sum_{i = 1}^{N_R} \nu_{i,j} \, r_i\]

Here, the summation goes over all reactions (\(i = 1\) to \(N_R\), where
\(N_R\) is the number of reactions occuring in the system). The
stoichiometric coefficient, \(\nu_{i,j}\), is the coefficient of the
\(j^{\text{th}}\) species in the \(i^{\text{th}}\) reaction, and \(r_i\)
is the rate of the \(i^{\text{th}}\) reaction. You'll notice that these
species production rates are still intensive and, for homogeneous
systems, have the following units, just as they did for a single
reaction:

\[R_j \ [=] \ \frac{\text{moles of }j}{\text{volume} \cdot \text{time}}\]

    \hypertarget{connecting-intensive-production-rates-r_j-to-extensive-production-rates-g_j}{%
\subsubsection{\texorpdfstring{Connecting Intensive production rates,
\(R_j\) to Extensive production rates,
\(G_j\)}{Connecting Intensive production rates, R\_j to Extensive production rates, G\_j}}\label{connecting-intensive-production-rates-r_j-to-extensive-production-rates-g_j}}

Recall that a species balance on a reactor is fundamentally an extensive
balance that changes with the size of the system and has units of
moles/time:

\[\frac{dN_j}{dt} = F_{j,f} - F_j + G_j\]

The \textbf{basic} idea of relating an intensive property of a system to
an extensive property of a system is that we simply scale the intensive
property to the size of the system.

\hypertarget{a-simple-example}{%
\paragraph{A Simple Example}\label{a-simple-example}}

Consider that we have a homogeneous reaction, \(A \longrightarrow B\),
that is occuring in a 1 Liter reactor. We will again call this
``Reaction 1'', and we will set the concentration of \(A\) at 1 mole per
liter, and there is no \(B\) initially in the system. At these
conditions, the intensive rate of Reaction 1 is:

\[r_1 = 1 \ \frac{\text{moles of extent of reaction} \ 1}{\text{L} \cdot \text{min}}\]

Based on our discussion above, the intensive production rates for A and
B would be:

\begin{align}
    R_A &= -1 \ \frac{\text{mol}}{\text{L} \cdot \text{min}}\\
    R_B &= 1 \ \frac{\text{mol}}{\text{L} \cdot \text{min}}\\
\end{align}

And to get the extensive production rates for A and B, we'd multiply
those intensive rates by the size of the system, which is 1L here:

\begin{align}
    G_A &= -1 \ \frac{\text{mol}}{\text{min}}\\
    G_B &= 1 \ \frac{\text{mol}}{\text{min}}\\
\end{align}

\hypertarget{scaling-to-100l}{%
\paragraph{Scaling to 100L}\label{scaling-to-100l}}

If we scale this system to a 100L reactor with identical temperature,
pressure, and composition, we find that it has an identical intensive
rate of reaction, which makes sense because reaction rate only depends
on \((T, P, x_j)\), which is identical in the two reactors:

\[r_1 = 1 \ \frac{\text{moles of extent of reaction} \ 1}{\text{L} \cdot \text{min}}\]

Similarly, intensive production rates are identical in the 100L reactor
and the 1L reactor:

\begin{align}
    R_A &= -1 \ \frac{\text{mol}}{\text{L} \cdot \text{min}}\\
    R_B &= 1 \ \frac{\text{mol}}{\text{L} \cdot \text{min}}\\
\end{align}

It is only when we get to calculating extensive rates--which we do by
multiplying intensive rates by the size of the system--that we see a
difference:

\begin{align}
    G_A &= -100 \ \frac{\text{mol}}{\text{min}}\\
    G_B &= 100 \ \frac{\text{mol}}{\text{min}}\\
\end{align}

\begin{figure}
\centering
\includegraphics{IntensiveExtensive.png}
\caption{Intensive vs.~Extensive Quantities in a Chemical Reactor}
\end{figure}

    \hypertarget{doing-this-for-a-real-reactor}{%
\subsubsection{Doing this for a real
reactor\ldots{}}\label{doing-this-for-a-real-reactor}}

What we did above, directly multiplying an intensive production rate by
the volume of the system is something we can only do if the intensive
reaction rate is \textbf{\emph{the same}} throughout the reactor volume.
This is not always true. Consider our arbitrary reactor:

\begin{figure}
\centering
\includegraphics{ArbitraryReactor.png}
\caption{Calculating Extensive Production Rates in an Arbitrary Reactor}
\end{figure}

This reactor has a total volume of \(V\). We can further conceptually
subdivide that volume into many small volumes, each with a volume of
\(V_k\). If we sum up all of those volumes, we would get the total
volume of the system:

\[V = \sum_k V_k\]

Now we consider potential differences between the small volume element
at position 1, \(V_1\), and the one at positions 2, \(V_2\), and 3,
\(V_3\). We know that rates of reaction are a function of only
temperature, pressure, and composition. So if we have a single reaction
\(i\) occuring in this reactor, we could define the reaction rate inside
of any volume element \(k\), \(r_{i,k}\) as:

\[r_{i,k} = f(T, P, x_j)_k\]

We have no way to guarantee that the temperature, pressure, and
compositions are equal at all of the positions in our reactor, so it is
entirely possible that:

\[(T, P, x_j)_1 \neq (T, P, x_j)_2 \neq (T, P, x_j)_3\]

Because \((T, P, x_j)_k\) fully determine the reaction rate within
volume element \(k\), it follows that it is possible for:

\[r_{i,1} \neq r_{i,2} \neq r_{i,3}\]

If this is true, then it also means that production rates may be
spatially non-uniform throughout the reactor:

\[R_{j,1} \neq R_{j,2} \neq R_{j,3}\]

It is important to note that we are still only considering a single
reaction--the subscripts 1, 2, and 3 refer to \textbf{\emph{positions}}
in the reactor. The implication is that, without making some specific
assumptions about our reactor, we \textbf{\emph{cannot}} guarantee that
reaction rates and production rates are spatially uniform everywhere
inside of the reactor. Since we can't guarantee spatially uniform rates,
we \textbf{\emph{cannot}} simply multiply the intensive production rate
of a species by the reactor volume to get the extensive production rate.
It's slightly more difficult\ldots{}

We start with our individual volume elements in the system, \(V_k\). We
will assume that each of these elements is small. Moreso, we will assume
that it is sufficiently small that we can say the temperature, pressure,
and composition \textbf{\emph{inside}} of a single element, \(V_k\) are
spatially uniform throughout that small volume element. If this is true,
then we can say that the intensive rates of all reactions are spatially
uniform throughout that small volume element. If this is true, then the
intensive production rate of each species must be spatially uniform
throughout that small volume element. If these things are true, then for
the small volume element, \(V_k\), we can calculate an extensive
production rate for species \(j\) in that volume element, \(G_{j,k}\) as
follows:

\[G_{j,k} = R_{j,k} V_k \ [=] \ \frac{\text{moles of} \ j}{\text{time}} \]

Since we know the extensive production rate in a single volume element,
\(V_k\), we can sum up the extensive rate in each element to calculate
the total extensive production rate of \(j\) throughout the entire
volume, \(V\).

\[G_{j} = \sum_{k = 1}^{m} G_{j,k} = \sum_{k = 1}^{m} R_{j,k} V_k \ [=] \ \frac{\text{moles of} \ j}{\text{time}} \]

Focusing on this part of the equation:

\[G_{j} = \sum_{k = 1}^{m} R_{j,k} V_k\]

If we assume that we have an infinite number of volume elements
(\(m \longrightarrow \infty\)), the volume of each element approaches
zero, and this summation becomes an integral:

\[G_{j} = \int^{V} R_j dV\]

In other words, we allow that \(R_j\) is a \textbf{\emph{function of
Volume}}, and we integrate it over the entire reactor volume to compute
the extensive production rate of species \(j\), \(G_j\). Without making
specific declarations about the type of reactor we are working with and
assumptions about how it operates, we cannot simplify this expression
any further. So, for our general species balance on a reactor:

\[\frac{dN_j}{dt} = F_{j,f} - F_j + G_j\]

The final result that we'll use (\textbf{\emph{for a homogeneous
reaction}}) is:

\[\frac{dN_j}{dt} = F_{j,f} - F_j + \int^{V} R_j dV\]

In the next lecture, we will look at how to apply this general balance
equation to various types of reactors (Batch, CSTR, PFR) operating under
specific assumptions.


    % Add a bibliography block to the postdoc
    
    
    
\end{document}
